% PROVENANCE (section-level)
%   produced-by : utilities.py :: create_templates_new(), create_workspace()
%   commit      : dd520a5
\section{Signal Extraction}
\label{sec:fit}

\subsection{Fit observables and channels}
\label{sec:fit:observables}
Yields are extracted from a two-dimensional binned likelihood fit in
$\mmiss$ (\texttt{B0\_recMissM2}) and $\pdl$ (\texttt{p\_D\_l}), performed
simultaneously in two channels:
\begin{itemize}
  \item the $D$-mass signal region, $1.855 < m(K\pi\pi) < 1.885$;
  \item the $D$-mass sideband, $[1.79,1.82] \cup [1.92,1.95]$
        (Sec.~\ref{sec:recon:d}), which constrains the fake-$D$ background.
\end{itemize}
No other control region enters the fit; the $q^{2}$, $D^{*}$-veto, BDT and
wrong-charge regions of Sec.~\ref{sec:validation} are cross-checks.
\AnalysisTBD{motivation for a 2D fit rather than two 1D fits, in terms of the
correlation between $\mmiss$ and $\pdl$ and the separation it buys}
\AnalysisTBD{the nominal binning in each channel, and the bin-merging and
coverage thresholds actually used}

\subsection{Templates and parameters}
\label{sec:fit:parameters}
Templates are built per truth category (Sec.~\ref{sec:truth}) with per-bin
statistical uncertainties from the sum of squared weights.
The parameter of interest is the normalisation of the $B \to D \tau \nu$
template; $R(D)$ is derived from it together with the normalisation of
$B \to D \ell \nu$, not fitted directly.
\AnalysisTBD{the explicit relation between the fitted normalisations and
$R(D)$, including the efficiency ratio and the $\tau \to \ell \nu \nu$
branching fraction}
\AnalysisTBD{full table of channels, samples, observables, bins, parameters
of interest, nuisance parameters, modifiers, constraints, initial values,
bounds and correlations, cross-checked against a generated pyhf workspace}

Two issues in the current workspace generator require a decision:
\AnalysisTBD{all normalisation parameters, including the signal, are given
bounds $[-5, 5]$, which permits negative yields}
\AnalysisTBD{the generic-$B\bar{B}$ families enter as free normalisation
factors.  This is appropriate for the tuning-script context but not for the
signal-region fit, where they should be fixed or constrained; otherwise the
final fit reopens what the tuning of Sec.~\ref{sec:bbbar} constrained}

\subsection{Fit validation}
\label{sec:fit:validation}
\AnalysisTBD{Asimov sensitivity, toy studies, pull distributions and
linearity tests, each reported separately from any result on data}
